\scanpage{374}
\mspage{113}
% archive sheet number: 373
on aurait une description du $\mathrm{Hot}$-groupe $\Omega(X,a)=\underline{\mathrm{Ch}}_{\unsure{\Sigma|Z}}(X,a)$ en termes d'une
% slanted
\margin{On fait ici $\Sigma=\unsure{\mathrm{Fl}|\mathrm{F}\ell}(X)$} catégorie $C$ munie d'une loi de composition \uline{associative} (et unitaire) \inserted{\uline{dans $\mathrm{Cat}$}}, \unsure{relevant|redonnant} la loi de groupe de $\Omega$ en passant de $(\mathrm{Cat})$ à $\mathrm{Hot}$. Prenant le Nerf de $C$, on trouverait un \unsure{monoïde|modèle} $\frac12$ simplicial qui, dans $\mathrm{Hot}$, donnerait la loi de groupe de $\Omega(X,a)$. Or je me rappelle \unsure{qu'on m'a|qu'il m'a} expliqué qu'en général, la \unsure{loi de groupe} de $\Omega(X,a)$ ne \unsure{saurait|pourrait} pas être \unsure{donnée|représentée} par une loi \uline{associative} sur un \unsure{ensemble} $\frac12$ simplicial. Ainsi, l'insuccès de mes efforts semble tenir, non à ma maladresse, mais à la nature des choses.
% slanted
\margin{C'est une \unsure{erreur}, -- ça doit être \unsure{plein|pour|faux} d'\unsure{``autres|autres} objets groupe'' de $\mathrm{Hot}$ que les groupes de la \unsure{catég.} \unsure{J'ai vu|j'ai lu} dans \unsure{Quillen} que $\Omega(X,x)$ doit bien se décrire par un \unsure{tel} groupe $\frac12$ simplicial (quand $X$ est un ens. q.c. $\frac12$ simplicial, à un sommet !). Donc \unsure{le problème|la question} si $C_\infty\simeq C$ dans $\mathrm{Hot}$ \unsure{reste} \unsure{entier|entière} !}

Je vais résumer les faits positifs auxquels je suis arrivé dans le diagramme suivant, qui a lieu dans la situation générale $C\mathrel{\unsure{-|\to}}T$ envisagée au début du présent paragraphe :
